\documentclass[11pt,a4paper]{article}
\usepackage[utf8]{inputenc}
\usepackage[T1]{fontenc}
\usepackage[margin=1in]{geometry}
\usepackage{amsmath}
\usepackage{amssymb}
\usepackage{amsfonts}
\usepackage{graphicx}
\usepackage{enumitem}
\usepackage{booktabs}
\usepackage{graphicx}
\usepackage{bm}

\newcommand{\sheetnumber}{3}
\newcommand{\R}{\mathbb{R}}

\usepackage{listings}

\providecommand{\xx}{\mathbf{x}}
\providecommand{\yy}{\mathbf{y}}
\providecommand{\XX}{\mathbf{X}}
\providecommand{\tth}{\boldsymbol{\theta}}
\providecommand{\sigm}{\sigma}
\providecommand{\rank}{\operatorname{rank}}

\lstset{
  basicstyle=\ttfamily\small,
  frame=single,
  breaklines=true,
  columns=fullflexible
}

\begin{document}

\begin{titlepage}
  \begin{center}
    {\Large \textbf{Exercise Sheet \sheetnumber}}

    \vspace{0.5em}
    Machine Learning for AIE, WS 2026/27

    \vspace{2em}
    {\LARGE \textbf{PCA and Model Selection: \\Predicting House Prices.}}
  \end{center}

  \vspace{2em}
  \begin{center}
    \begin{tabular}{|l|l|l|}\hline
      \multicolumn{3}{|c|}{\textbf{Team Members}} \\\hline
      Last name & First name & Matriculation number \\\hline
      ZENZ & Sebastian & 01031159 \\\hline
    \end{tabular}
  \end{center}
\end{titlepage}

\tableofcontents
\clearpage

\section{Task 1: Principal Component Analysis }
\label{sec:PCA}
 
\subsection{Task 1.1: Pen and paper computation}
\label{subsec:PCA1}

\begin{figure}[h]
  \centering
  \includegraphics[width=0.75\linewidth]{plots/Ex1_1_1.png}
\end{figure}

\begin{figure}[h]
  \centering
  \includegraphics[width=0.75\linewidth]{plots/Ex1_1_2.png}
\end{figure}

\clearpage
\subsection{Task 1.2: Implementation of PCA}
\label{subsec:PCA2}
The function \texttt{standardize\_and\_fit\_pca} was implemented and a PCA was fitted to the 9 features.
The eigenvalues and explained variance ration from the fitted PCA was computed and are listed in the table below.
With the inverse transformation provided in the \texttt{Scikit-Learn} module, the reconstruction error with the first 
K principal directions as well:

\begin{center}
\small
\begin{tabular}{rrrrr}
\toprule
$K$ & Eigenvalue $\lambda_K$ & Ratio & Cumulative ratio & Reconstruction MSE\\
\midrule
1 & 3.677 & 0.408 & 0.408 & 5.324\\
2 & 1.391 & 0.155 & 0.563 & 3.933\\
3 & 1.068 & 0.119 & 0.682 & 2.866\\
4 & 0.773 & 0.086 & 0.767 & 2.093\\
5 & 0.603 & 0.067 & 0.834 & 1.490\\
6 & 0.534 & 0.059 & 0.894 & 0.956\\
7 & 0.409 & 0.045 & 0.939 & 0.548\\
8 & 0.359 & 0.040 & 0.979 & 0.189\\
9 & 0.189 & 0.021 & 1.000 & 0.000\\
\bottomrule
\end{tabular}
\end{center}

\begin{enumerate}[label=\textbf{Q1.\arabic*:}]
\item Standardization ist important here because the features have very different units and ranges (see output of variance of the raw features below).
Without standardization, lot area could dominate because its raw
variance is much larger than the other features' variances

\begin{lstlisting}
variance of the raw features:
  living_area          2013.1
  lot_area           583330.6
  basement_area        1443.1
  garage_area           350.9
  year_built            869.2
  overall_quality         1.8
  bedrooms                0.7
  bathrooms               0.3
  fireplaces              0.4
\end{lstlisting}

\item The reconstruction error using the first $K$ directions is approximately
the sum of the omitted eigenvalues (as seen in in Task 1.1(b) as well).
For compression, one can retain the smallest $K$ that still meets a
defined variance target: here $K=7$ retains $93.9\%$ of variance, while $K=6$
retains $89.4\%$.
\\The Figure~\ref{fig:expVar} plots the trend between (cumulative) explained error and mean squared reconstruction error
in relation to the number of principal directions K used for the PCA.

\begin{figure}[h]
  \centering
  \includegraphics[width=0.75\linewidth]{plots/1_explained_variance.pdf}
  \caption{Explained variance and MSCE over K}
  \label{fig:expVar}
\end{figure}

\item The first principal directions of b\_1 are included in the output below. 
All PC1 scores are positive. The largest are living area ($0.42$),
overall quality ($0.41$), bathrooms ($0.39$), garage area ($0.37$), and
basement area ($0.34$). PC1 therefore summarizes a general combination of
house size, quality, etc. A larger $z_1$ indicates a larger,
better-equipped house (like a new pseudo feature).
\begin{lstlisting}
first principal direction b_1:
  living_area        0.42
  lot_area           0.16
  basement_area      0.34
  garage_area        0.37
  year_built         0.31
  overall_quality    0.41
  bedrooms           0.20
  bathrooms          0.39
  fireplaces         0.29
\end{lstlisting}

\item Although PCA does not use prices as an input, these housing characteristics are
associated with price: larger, newer, higher-quality houses tend to sell for
more. Thus the feature-based scores can correlate with price, without implying
that PCA learned or causes that relationship. 
\\Figure~\ref{fig:pcScatter} shows the housing price in relation to the scores $z_1$ and $z_2$. The positive correlation
of increasing house prices with increasing magnitude of the feature scores (especially $z_1$) is clearly visible.

\begin{figure}[h]
  \centering
  \includegraphics[width=1\linewidth]{plots/1_pc_scatter.pdf}
  \caption{Scatter plot of house price over PCA scores z1, z2}
  \label{fig:pcScatter}
\end{figure}
\end{enumerate}

\clearpage
\section{Task 2: Overfitting and Model Evaluation}
\label{sec:overfit}

\subsection{Task 2.1: Toy example}

For degree $1$, the fitted parameters are $\theta=(0.5,0.6)$, the training
MSE is $0.050$, and the squared error on the fifth house is $0.010$.
For degree $3$, $\theta=(-3,6.167,-2.5,0.333)$, the training MSE is
approximately zero, and the test-house squared error is $12.960$

\begin{lstlisting}
  d = 1: theta = [0.5 0.6], training MSE = 0.050, squared error on the test house = 0.010
  d = 3: theta = [-3.     6.167 -2.5    0.333], training MSE = 0.000, squared error on the test house = 12.960
  design matrix Phi of degree 3:
  [[ 1.  1.  1.  1.]
  [ 1.  2.  4.  8.]
  [ 1.  3.  9. 27.]
  [ 1.  4. 16. 64.]]
\end{lstlisting}

\begin{enumerate}[label=\textbf{Q2.\arabic*:},start=1]
\item Here, the design matrix is 4x4: four training points and four cubic-model parameters. 
The rank of the matrix is 4 as well (full-rank). Since the four input values are
different, a polynomial with degree 3 can fit all four training points exactly and therefore the training MSE equals to 0.
\item  The cubic model achieves the lower training loss, if that would be our only criterion on choosing the model complexity,
the polynomial degree of 3 would be the better choice. Choosing the cubic leads to a higher test loss, as the higher model
complexity results in a more flexibel model, that fits the training data better, but has a higher test loss or: it generalizes
worse for unseen data. This phenomenon is called overfitting!
\end{enumerate}

\subsection{Task 2.2: Training, test, and holdout validation}

With dataset was split into 100 houses for the training set and 2313 test houses. The polynomial degree is treated as a hyperparameter
for this task and the training and test MSE for d=1..6 are reported:
\begin{lstlisting}
100 training houses, 2313 test houses
  d = 1: training MSE = 780.4, test MSE = 1162.9
  d = 2: training MSE = 657.9, test MSE = 904.8
  d = 3: training MSE = 616.6, test MSE = 903.6
  d = 4: training MSE = 585.0, test MSE = 1962.3
  d = 5: training MSE = 582.3, test MSE = 4214.3
  d = 6: training MSE = 574.1, test MSE = 14568.3
\end{lstlisting}
\vspace{1em}
\noindent Figure~\ref{fig:mvsPolyFits} plots the polynomials of degree 1, 2 and 6 for the fitted model. With the the polynomial degree 6 the
the overfitting phenomena is illustrated pretty well. The model is flexible enough, to find the optimum to approximate to the training
dataset (black dots). As the training dataset is rather small, the model tends to overfit to the seen training data but fails
to generalized the underlying trend of the complete dataset (all dots, including opaque grey datapoints) very well.

\clearpage
\begin{figure}[h]
  \centering
  \includegraphics[width=1\linewidth]{plots/2_polynomial_fits.pdf}
  \caption{Polynomial fits over first principle direction z1}
  \label{fig:mvsPolyFits}
\end{figure}


\noindent For Task 2.2(b) the holdout method is used. The training set is further split into a (sub-)training and a validation set.
The model is fitted and trained on the training set and validated with the validation set.
Bases ond the lowest validation loss (d=3 in this case), the model is selected and the validation and subtraining set are combined again for
a final training cycle and evaluated against the test set.


\begin{lstlisting}
  d = 1: validation MSE = 893.9
  d = 2: validation MSE = 1210.0
  d = 3: validation MSE = 799.2
  d = 4: validation MSE = 2770.7
  d = 5: validation MSE = 3403.9
  d = 6: validation MSE = 2562.0
  selected degree: d = 3, test MSE after re-training on all 100 training houses = 903.6
\end{lstlisting}

\noindent The holdout method selects degree $3$; after fitting it on all 100 training
houses, its test MSE is $903.6$.
\clearpage
\noindent Figure~\ref{fig:mvsErrors} plots the calculated training, validation and test MSE over the polynomial degree d.
In this figure the effect of overfitting can also be detected. The training loss is getting lower with increasing polynomial degree
(models fits the training set better with higher complexity), but the training and validation error is increasing --> model generalizes
worse for unseen data. Selecting a model only with the training loss will result in an overly optimistic estimate for the real
test/use cases and will result in bad performance.

\begin{figure}[h]
  \centering
  \includegraphics[width=1\linewidth]{plots/2_errors.pdf}
  \caption{Training, validation and test MSE over polynomial degree }
  \label{fig:mvsErrors}
\end{figure}

\begin{enumerate}[label=\textbf{Q2.\arabic*:},start=3]
\item Polynomial model classes are nested: a degree-$d+1$ polynomial can
represent any degree-$d$ polynomial by setting a coefficient to zero.
Therefore, the minimum training error cannot increase with degree.
\\ The model with polynomial degree $1$ underfits (high training and test error) relative to the better-performing degrees $2$ and $3$.
Degrees $4$--$6$ overfit: their training errors continue to decrease while
their test errors increase. 
\\Figure~\ref{fig:mvsPolyFits} indicates, that the largest test error are meassured
at the extreme principal scores (far from the mean), while estimating the dense "middle part", close to the mean, pretty well.
\item The test set MUST not be used to select the degree because selection
adapts the model to test outcomes (feedback) and makes the reported test error optimistically
biased. The test MSE of the selected model is no longer an independent,
unbiased final estimate.
\item A validation set of 30 houses gives a very noisy estimate, so a different
random split may change the calculated losses significantly and suggest another hyperparameter configuration. 
The selected degree should therefore not be treated as certain; cross-validation can reduce dependence on one
particular split.
\end{enumerate}

\clearpage
\section{Task 3: Model Selection with Cross-Validation}
The five-fold validation results and their mean are:
\begin{center}
\small
\begin{tabular}{rrrrrrr}
\toprule
$d$ & Fold 1 & Fold 2 & Fold 3 & Fold 4 & Fold 5 & CV MSE\\
\midrule
1 & 847.8 & 1329.5 & 875.8 & 487.6 & 563.4 & 820.8\\
2 & 920.6 & 1084.6 & 686.0 & 422.0 & 519.6 & 726.6\\
3 & 1141.1 & 969.1 & 646.2 & 421.4 & 451.1 & 725.8\\
4 & 1607.7 & 871.4 & 620.2 & 405.1 & 372.6 & 775.4\\
5 & 2115.1 & 952.5 & 613.7 & 415.0 & 365.6 & 892.4\\
6 & 2098.5 & 971.8 & 600.6 & 417.4 & 354.5 & 888.6\\
\bottomrule
\end{tabular}
\end{center}
Cross-validation selects $d=3$ (CV MSE $725.8$). Retraining on all
100 training houses gives test MSE $903.6$.

\begin{figure}[h]
  \centering
  \includegraphics[width=0.75\linewidth]{plots/3_cross_validation.pdf}
  \caption{Cross validation and test MSE over polynomial degrees}
  \label{fig:cvMSE}
\end{figure}

\begin{enumerate}[label=\textbf{Q3.\arabic*:}]
\item If each fold is used as the validation set separately, the selected
degrees are: fold 1: $d=1$; fold 2: $d=4$; fold 3: $d=6$; fold 4: $d=5$;
fold 5: $d=6$. These vary substantially, whereas averaging across folds
selects $d=3$.
\item The holdout method and cross-validation both select $d=3$ in this run.
Cross-validation is generally more trustworthy here because it averages over
five validation splits instead of relying on one 30-house set. The mean CV
errors for degrees 2 and 3 are close, however, so the selection is not
certain.
\item Using identical folds makes the degree comparison paired: each degree
is evaluated on the same held-out examples. This reduces variation due to
different splits and makes differences in validation error more attributable
to the degree.
\item Holdout is preferable when the dataset is large enough for a
representative validation set and model training is costly. It requires fewer
fits than $K$-fold CV, while CV is useful when data is limited and a more
stable estimate is worth the additional computation.
\end{enumerate}

\clearpage
\section{Addendum}
\section{Usage of generative AI}
Generative AI was used in the the form of GitHub Copilot in Visual Studio Code. As the coding effort for this assignment
was managable, AI was mostly used to get information on how to properly use the different scikit-learn function, without
the need to search the documentation on the scikit website and saving time this way.
\\Another usecase was the import of tables (pandas etc.) from visual studio to a nicely formatted table in LaTeX.

\end{document}